\chapter{กฎการเคลื่อนที่ของนิวตันและแรงกระทำ}
\label{ch:ch04_newton_laws}
\index{กฎการเคลื่อนที่ของนิวตัน}

\begin{conceptbox}[title=วัตถุประสงค์การเรียนรู้ประจำบท]
เมื่อสิ้นสุดการศึกษาบทเรียนนี้ นักศึกษาสามารถ
\begin{enumerate}
    \item เมื่อศึกษาบทนี้จบแล้ว ผู้เรียนควรมีความสามารถดังต่อไปนี้
    \item อธิบายกฎการเคลื่อนที่ทั้งสามข้อของเซอร์ไอแซก นิวตัน พร้อมยกตัวอย่างประกอบจากสถานการณ์จริง
    \item เขียนแผนภาพวัตถุอิสระ (free-body diagram) ของวัตถุในสถานการณ์ต่าง ๆ ได้อย่างถูกต้อง
    \item ประยุกต์สมการ $\Sigma$F = ma เพื่อวิเคราะห์และแก้โจทย์ปัญหาการเคลื่อนที่ในหนึ่งและสองมิติ
    \item อธิบายความแตกต่างระหว่างแรงเสียดทานสถิตและแรงเสียดทานจลน์ พร้อมคำนวณค่าแรงเสียดทานในสถานการณ์ต่าง ๆ
\end{enumerate}
\end{conceptbox}

\section{ทฤษฎีและหลักการพื้นฐาน}

\begin{bioappbox}
\textbf{แรงหนืดของสโตกส์และกลศาสตร์การเคลื่อนที่ในของไหลชีวภาพ}

เมื่อจุลินทรีย์เคลื่อนที่ในน้ำหรือสารคัดหลั่ง แรงเสียดทานหลักที่กระทำต่อเซลล์คือแรงต้านความหนืดตามกฎของสโตกส์ (Stokes' Law) $F_d = 6\pi\eta r v$ เนื่องจากขนาดอนุภาคมีค่าน้อยมาก ค่าเรย์โนลด์สนัมเบอร์ (Reynolds number) ของแบคทีเรียจึงมีค่าน้อยกว่า $10^{-4}$ ซึ่งหมายความว่าแรงหนืดมีความสำคัญเหนือกว่าแรงเฉื่อยอย่างสิ้นเชิง สิ่งมีชีวิตในระดับจุลภาคจึงต้องใช้แฟลกเจลลา (Flagella) หมุนควงแบบเกลียวเพื่อขับเคลื่อนตัวเองไปข้างหน้า
\end{bioappbox}

\begin{labbox}
1. ตรวจสอบตำแหน่งศูนย์ (Zero Error) ของเครื่องมือวัดทุกชนิดก่อนเริ่มบันทึกข้อมูล
2. ทำการวัดซ้ำอย่างน้อย 3 ถึง 5 ครั้งในแต่ละจุดทดลองเพื่อคำนวณหาค่าเฉลี่ยและค่าเบี่ยงเบนมาตรฐาน
3. บันทึกผลการวัดพร้อมระบุหน่วยเอสไอ (SI Units) และค่าความคลาดเคลื่อนของการวัดเสมอ
\end{labbox}

\section{ตัวอย่างโจทย์และการวิเคราะห์เชิงฟิสิกส์}

\begin{examplebox}
ตัวอย่างที่พบเห็นได้ในชีวิตประจำวัน เช่น เมื่อรถยนต์เบรกกะทันหัน ผู้โดยสารที่ไม่ได้คาดเข็มขัดนิรภัยจะพุ่งไปข้างหน้า เนื่องจากร่างกายของผู้โดยสารยังคงมีแนวโน้มรักษาสภาพการเคลื่อนที่เดิมไว้ตามกฎความเฉื่อย ในทำนองเดียวกัน เมื่อดึงผ้าปูโต๊ะออกอย่างรวดเร็วจากใต้จานที่วางอยู่ จานจะยังคงอยู่ในตำแหน่งเดิมได้ชั่วขณะเนื่องจากแรงเสียดทานระหว่างผ้ากับจานมีเวลากระทำน้อยมาก กฎข้อที่หนึ่งจึงเป็นการนิยามกรอบอ้างอิงเฉื่อย (inertial reference frame) ซึ่งเป็นกรอบอ้างอิงที่กฎการเคลื่อนที่ของนิวตันใช้ได้อย่างสมบูรณ์

4.2 กฎข้อที่สอง: $\Sigma$F = ma

กฎการเคลื่อนที่ข้อที่สองของนิวตันเป็นหัวใจเชิงปริมาณของกลศาสตร์คลาสสิก โดยกล่าวว่าความเร่งของวัตถุแปรผันตรงกับแรงลัพธ์ที่กระทำต่อวัตถุ และแปรผกผันกับมวลของวัตถุ เขียนเป็นสมการได้ว่า

$\Sigma$F = ma
\end{examplebox}

\section{แบบฝึกหัดทบทวนและคำถามท้ายบท}

\begin{enumerate}
    \item จงอธิบายกฎการเคลื่อนที่ทั้งสามข้อของนิวตันด้วยคำพูดของตนเอง พร้อมยกตัวอย่างจากชีวิตประจำวันข้อละ 1 ตัวอย่าง
    \item เพราะเหตุใดวัตถุที่มีมวลมากจึงมีความเฉื่อยมากกว่าวัตถุที่มีมวลน้อย จงอธิบาย
    \item กล่องมวล 12 kg วางบนพื้นราบไม่มีความเสียดทาน ถูกผลักด้วยแรงในแนวราบขนาด 36 N จงหาความเร่งของกล่อง
    \item วัตถุมวล 8 kg ถูกดึงด้วยแรง 60 N ในแนวราบ พบว่าวัตถุมีความเร่ง 4 m/s$^2$ จงหาแรงเสียดทานที่กระทำต่อวัตถุ
    \item กล่องมวล 25 kg วางบนพื้นราบ มีสัมประสิทธิ์แรงเสียดทานจลน์ $\mu$\_k = 0.30 จงหาแรงเสียดทานจลน์ที่กระทำต่อกล่องขณะเคลื่อนที่
    \item วัตถุมวล 5 kg วางนิ่งบนพื้นราบ มี $\mu$\_s = 0.50 จงหาแรงเสียดทานสถิตสูงสุดที่พื้นสามารถกระทำต่อวัตถุได้
\end{enumerate}

% สิ้นสุดบทเรียน
